\documentclass[12pt,a4paper]{article}
\usepackage[utf8]{vntex, inputenc}
\usepackage{amsmath, amssymb, anyfontsize, array, color,  enumerate, enumitem, fancyhdr, float, graphicx, hyperref, lastpage, listings, mathtools, mips, natbib, pifont, tikz}
\usepackage[open,openlevel=0]{bookmark}
\usepackage[framemethod=tikz]{mdframed}

\usetikzlibrary{calc}
\newcommand\HRule{\rule{12cm}{1pt}}

\pagestyle{fancy}
\fancyhead{} % clear all header fields
\fancyhead[L]{
  \color{blue}
  \begin{tabular}{rl}
    \begin{picture}(25,15)(0,0)
    \put(0,-8){\includegraphics[width=8mm, height=8mm]{BK.png}}
    \end{picture}
    \begin{tabular}{l}
      \textbf{\bf \ttfamily Ho Chi Minh University of Technology}\\
      \textbf{\bf \ttfamily Faculty of Computer Science \& Engineering}
    \end{tabular}
  \end{tabular}
}
\fancyhead[R]{
  \begin{tabular}{l}
    \tiny \bf \\
    \tiny \bf
  \end{tabular}
}
\fancyfoot{} % clear all footer fields
\fancyfoot[L]{\scriptsize \ttfamily Computer Architecture Week 5}
\fancyfoot[R]{\scriptsize \ttfamily Page {\thepage}/\pageref{LastPage}}
\renewcommand{\headrulewidth}{0.3pt}
\renewcommand{\footrulewidth}{0.3pt}

\definecolor{dkgreen}{rgb}{0,0.6,0}
\definecolor{gray}{rgb}{0.5,0.5,0.5}
\definecolor{mauve}{rgb}{0.58,0,0.82}

\lstset{frame=tb,
  language=[MIPS]Assembler,
  aboveskip=3mm,
  belowskip=3mm,
  showstringspaces=false,
  columns=flexible,
  basicstyle={\small\ttfamily},
  numbers=none,
  numberstyle=\tiny\color{gray},
  keywordstyle=\color{blue},
  commentstyle=\color{dkgreen},
  stringstyle=\color{mauve},
  breaklines=true,
  breakatwhitespace=true,
  tabsize=3,
  numbers=left,
  stepnumber=1,
  numbersep=1pt,
  firstnumber=1,
  numberfirstline=true
}

\begin{document}
\begin{titlepage}
  \begin{tikzpicture}[remember picture, overlay]
    \draw[line width = 2pt] ($(current page.north west) + (1in,-1in)$) rectangle ($(current page.south east) + (-1in,1in)$);
  \end{tikzpicture}

  \begin{center}

    % Upper part of the page
    \textbf{\fontsize{12pt}{1pt}\selectfont HO CHI MINH CITY UNIVERSITY OF TECHNOLOGY}\\[0.5cm]
    {\fontsize{13pt}{1pt}\selectfont Faculty of Computer Science \& Engineering}\\[0.5cm]
    \begin{figure}[H]
      \centering
      \includegraphics[width=1.7in,height=1.7in]{BK.png}
    \end{figure}

    % Title
    \HRule\\[0.5cm]
    { \textbf{\fontsize{25pt}{1pt}\selectfont Computer Architecture Week 5}}\\[0.4cm]

    \HRule\\[0.8cm]
    \begin{minipage}{0.545\textwidth}
      \begin{flushleft}
        \textbf{Authors:}\\
        \begin{tabular}{l l}
          Lưu Nguyễn Hoàng Minh & 1952845 \\
          Vũ Anh Nhi            & 1952380 \\
          Nguyễn Chính Khôi     & 1952793 \\
          Nguyễn Hoàng          & 1952255 \\
        \end{tabular}
      \end{flushleft}
    \end{minipage}
    \begin{minipage}{0.4\textwidth}
      \begin{flushright}
        \textbf{Class:}\\
        CO2008\_CC02\\
        \textbf{Lecturer:}\\
        Kiều Đỗ Nguyên Bình\\

      \end{flushright}
    \end{minipage}

    \vfill

    % Bottom of the page
    \vspace{2cm}
    {\large} %{\large \today}
  \end{center}
\end{titlepage}


\pdfbookmark[0]{Question 1}{ques1}
\section*{Question 1:}
Students run the program and answer the following question

\begin{mdframed}[hidealllines=true,backgroundcolor=magenta!10]
  \begin{lstlisting}

    .data:
    endl: .asciiz "\n"

    # Start program code
    .text:

    .globl main
  main:

    ## MULTIPLICATION
    lui $t0, 0x5000
    li $t1, 4
    mult $t0, $t1
    mflo $s0
    mfhi $s1

    # Print $t0
    li      $v0, 1
    move    $a0, $t0
    syscall

    # Print new line
    li      $v0, 4
    la      $a0, endl
    syscall

    # Print $t0
    li      $v0, 1
    move    $a0, $t1
    syscall

    # Print new line
    li      $v0, 4
    la      $a0, endl
    syscall

    # Print $s0
    li      $v0, 1
    move    $a0, $s0
    syscall

    # Print new line
    li      $v0, 4
    la      $a0, endl
    syscall

    # Print $s1
    li      $v0, 1
    move    $a0, $s1
    syscall

    # Print new line
    li      $v0, 4
    la      $a0, endl
    syscall


    ## DIVISION
    li $t0, 10
    li $t1, 7
    div $t0, $t1
    mflo $s0
    mfhi $s1

    # Print $s0
    li      $v0, 1
    move    $a0, $s0
    syscall

    # Print new line
    li      $v0, 4
    la      $a0, endl
    syscall

    # Print Ss1
    li      $v0, 1
    move    $a0, $s1
    syscall


    ## EXIT
    li $v0, 10
    syscall

  \end{lstlisting}
\end{mdframed}

\begin{enumerate}
  \item What is the function of mflo and mfhi instruction?

        For multiplication:\\
        - mflo stores the least-significant 32 bits of the product of the multiplication into the lo register. \\
        - mfhi stores the most-significant 32 bits of the product of the multiplication into the hi register. \\

        For division:\\
        - mflo stores the quotient of the division operation into the lo register. \\
        - mfhi stores the remainder of the division operation into the hi register.\\

  \item What is the value of $\$s0$ and $\$s1$ after the multiplication.

        - Value of $\$s0$ : 0x40000000 \\
        - Value of $\$s1$ : 0x00000001 \\

  \item What is the value of $\$s0$ and $\$s1$ after the division?

        - Value of $\$s0$ : 1 \\
        - Value of $\$s1$ : 3 \\

  \item What is the role of lo and hi register in multiplication?

        - lo: least-significant 32 bits of the 64-bit product.\\
        - hi: most-significant 32 bits of the 64-bit product. \\

  \item What is the role of lo and hi register in division?

        - lo: 32-bit quotient\\
        - hi: 32-bit remainder \\

  \item Do hi and lo is a part of 32 general purpose register of MIPS processors?

        - No they are not a part of the 32 general purpose registers.

\end{enumerate}

\section*{Question 2:} Given the following log 2 procedure in C:
\begin{figure}[h]
  \centering
  \includegraphics[scale=1.4]{log2.PNG}
\end{figure}

Assume that all input is the exponent of 2. Implement a MIPS program that:
\begin{enumerate}
  \item Receive input value from user.
  \item Call the log 2 procedure.
  \item Get the return value of log 2 procedure and print to screen
\end{enumerate}

\begin{mdframed}[hidealllines=true,backgroundcolor=magenta!10]
  \begin{lstlisting}

    .data
    msg: .asciiz "Enter an integer: "
    result: .asciiz "Result: "

    .text:

    .globl main

  main:
    li $v0, 4
    la $a0, msg
    syscall

    li $v0, 5
    syscall
    move $a0, $v0

    jal log_2
    move $s0, $v0

    li $v0, 4
    la $a0, result
    syscall

    li $v0, 1
    move $a0, $s0
    syscall

    li $v0, 10
    syscall

  log_2:
    addi $t0, $t0, 0


  for:
    div $a0, $a0, 2
    beq $a0, 0, exit
    addi $t0, $t0, 1
    j for
  exit:
    move $v0, $t0
    jr $ra

  \end{lstlisting}
\end{mdframed}


\section*{Question 3:}Implement the following C code by using MIPS code. Assume that b and c are 10 and 7, respectively while input variable is read from keyboard. Print value of a to the terminal.

\begin{figure}[h]
  \centering
  \includegraphics[scale=1.4]{switch.PNG}
\end{figure}

\begin{mdframed}[hidealllines=true,backgroundcolor=magenta!10]
  \begin{lstlisting}
    .data:
    msg: .asciiz "Enter an integer: "
    print: .asciiz "a = "
    printlo: .asciiz "\nlo = "
    printhi: .asciiz "\nhi = "
    break: .asciiz "Please input other integer numbers"

    # Start program code
    .text:

    .globl main
  main:
    addi $s1, $zero, 10         # b = 10
    addi $s2, $zero, 7          # c = 7

    # Print string msg
    li      $v0, 4
    la      $a0, msg
    syscall

    # Get input and save
    li      $v0, 5
    syscall
    move    $s0, $v0

    add $t0, $s1, $s2           # b + c
    sub $t1, $s1, $s2           # b - c
    mult $s1, $s2               # b * c
    mflo $t2
    mfhi $t3
    div $s1, $s2                # b / c
    mflo $t4
    mfhi $t5            # b % c

    beq $s0, 0, case0
    beq $s0, 1, case1
    beq $s0, 2, case2
    beq $s0, 3, case3
    beq $s0, 4, case4
    j default

  case0:
    # Print string print
    li      $v0, 4
    la      $a0, print
    syscall

    # Print a
    li      $v0, 1
    move    $a0, $t0
    syscall

    j Exit

  case1:
    # Print string print
    li      $v0, 4
    la      $a0, print
    syscall

    # Print a
    li      $v0, 1
    move    $a0, $t1
    syscall

    j Exit

  case2:
    # Print string print
    li      $v0, 4
    la      $a0, print
    syscall

    # Print a
    li      $v0, 1
    move    $a0, $t3
    syscall
    # Print a
    li      $v0, 1
    move    $a0, $t2
    syscall

    # Print string print
    li      $v0, 4
    la      $a0, printlo
    syscall

    # Print a
    li      $v0, 1
    move    $a0, $t2
    syscall

    # Print string print
    li      $v0, 4
    la      $a0, printhi
    syscall

    # Print a
    li      $v0, 1
    move    $a0, $t3
    syscall

    j Exit

  case3:
    # Print string print
    li      $v0, 4
    la      $a0, print
    syscall

    # Print a
    li      $v0, 1
    move    $a0, $t4
    syscall

    j Exit

  case4:
    # Print string print
    li      $v0, 4
    la      $a0, print
    syscall

    # Print a
    li      $v0, 1
    move    $a0, $t5
    syscall

    j Exit

  default:
    # Print string break
    li      $v0, 4
    la      $a0, break
    syscall

  Exit:
    # Exit
    li      $v0, 10
    syscall

  \end{lstlisting}
\end{mdframed}
\end{document}
